\section{Discussion}
\label{sec:discussion}

\paragraph{Why does augmentation eliminate double descent?}
Our results show that all four augmentation strategies eliminate the double descent peak in ResNets, despite operating through quite different mechanisms---which itself is a somewhat surprising finding. We propose three complementary explanations, though we suspect the first is doing the most work.

\textit{(1) Effective sample size increase.} Double descent occurs near the interpolation threshold---the point where model capacity approximately equals the number of effective training constraints. Augmentation strategies that generate genuinely novel training examples (Cutout masking different patches per epoch; MixUp creating convex combinations unseen in the original data) increase the effective number of distinct training constraints the model must satisfy. This raises the interpolation threshold, potentially pushing it beyond our largest tested model size (2.78M), which would explain why the second descent never appears under augmentation. We find this explanation most convincing because it accounts for why even simple Flip+Crop---which adds no new semantic content---is sufficient to eliminate the peak: the additional geometric variants still expand the effective constraint set.

\textit{(2) Loss landscape smoothing.} MixUp is known to flatten decision boundaries by training on convex combinations of examples~\cite{zhang2018mixup}. Cutout encourages distributed representations by masking local patches, reducing reliance on texture cues and inducing more globally coherent feature representations. Both effects smooth the loss surface near the interpolation threshold, reducing the severity of the memorization-dominated regime that produces the peak. The well-established link between flat minima and generalization~\cite{foret2021sharpness,li2018visualizing} supports this account, though it is less clear why landscape smoothing would fully eliminate the threshold rather than merely reducing the peak height.

\textit{(3) Regularization-induced inductive bias shift.} Geometric augmentations (Flip+Crop) introduce equivariance priors, while Color Jitter reduces the model's reliance on color statistics. These inductive bias changes may reduce effective model complexity~\cite{nakkiran2020deep}, making the model behave as if it has fewer free parameters---effectively shifting the interpolation threshold rightward. This explanation feels less satisfying to us than (1), since it is harder to quantify and does not obviously generalize to MixUp, which does not introduce spatial priors.

\paragraph{Why is MixUp $\alpha$ insensitive?}
Honestly, we did not expect the five $\alpha$ curves to be this flat---the values are identical to four decimal places, not just approximately similar. Prior work~\cite{zhang2018mixup} reports that $\alpha$ affects downstream performance on standard benchmarks, so the near-zero sensitivity we observe is a genuine surprise. Our best interpretation: any nonzero $\alpha$ creates a training distribution that the model cannot perfectly interpolate (since mixed examples have fractional labels that induce soft constraints), which is sufficient to push the threshold past our experimental range. The degree of mixing modulates the ``softness'' of constraints but not their count, and the threshold appears more sensitive to the latter. That said, we cannot rule out that this flatness is specific to MLP on CIFAR-10. Because our ablation was conducted on MLP---an architecture that does not exhibit double descent---we cannot directly conclude that $\alpha$ is insensitive in the context of double descent suppression on ResNet. The same ablation on ResNet remains untested and may show different sensitivity, which we leave as an open question.

\paragraph{Why do MLPs not show double descent?}
The absence of double descent in MLPs---regardless of augmentation---was one of the more striking results we observed, since we expected both architectures to show the phenomenon at some capacity scale. We offer two hypotheses. First, the convolutional inductive bias in ResNets enforces local feature detectors whose capacity scales differently with width than global fully-connected features; this scaling mismatch may create sharper interpolation thresholds. Second, batch normalization (present in ResNet, absent in MLP) is known to interact with the interpolation regime~\cite{nakkiran2020deep} and may be a contributing factor. Our intuition leans toward batch normalization as the more proximate cause---removing it from ResNet while keeping the convolutional structure would be a cleaner test---but we leave this disentanglement to future work.

\paragraph{Limitations.}
Our experiments are conducted on CIFAR-scale datasets ($32{\times}32$ images, $\leq$50K training samples) and models with at most 2.78M parameters. Whether the same mechanisms operate in large-scale regimes---where the interpolation threshold may lie at billions of parameters---is an open question. Additionally, we study augmentation strategies independently; combinations may interact non-additively (e.g., MixUp + Cutout may not simply sum their individual effects). Finally, our mechanistic account remains explanatory rather than formal; a theoretical treatment connecting augmentation statistics to interpolation threshold position would strengthen these claims.
